\documentclass[a4paper,11pt]{ltjsarticle}
% Shared preamble for the proof documents in docs/proofs/.
% Build with LuaLaTeX:  lualatex level2-completeness.tex  (run twice for references)
\usepackage{amsmath,amssymb,amsthm}
\usepackage{mathtools}
\usepackage{booktabs}
\usepackage{longtable}
\usepackage{xcolor}
\usepackage[hidelinks]{hyperref}

\theoremstyle{plain}
\newtheorem{theorem}{定理}[section]
\newtheorem{lemma}[theorem]{補題}
\newtheorem{proposition}[theorem]{命題}
\newtheorem{corollary}[theorem]{系}
\theoremstyle{definition}
\newtheorem{definition}[theorem]{定義}
\newtheorem{check}{検査}
\renewcommand{\thecheck}{C\arabic{check}}
\theoremstyle{remark}
\newtheorem{remark}[theorem]{注意}

\newcommand{\Sym}{\operatorname{Sym}}
\newcommand{\Alt}{\operatorname{Alt}}
\newcommand{\sgn}{\operatorname{sgn}}
\newcommand{\supp}{\operatorname{supp}}
\newcommand{\Span}{\operatorname{span}}
\newcommand{\F}{\mathbf{F}}
\newcommand{\Z}{\mathbf{Z}}
\newcommand{\code}[1]{\texttt{\detokenize{#1}}}


\title{メンガー・ツイスト・キューブ Level 3 の\\構成的可解性}
\author{Menger Cube プロジェクト}
\date{2026年10月4日}

\begin{document}
\maketitle

\begin{abstract}
Level 3 のメンガー・ツイスト・キューブ（8,000セル）について，現行の衝突判定を通る3,591個の合法原子手から到達可能な任意の状態を，
実装 \code{level3-slice-reduction} が有限個の合法原子手で完成状態へ戻すことを証明する．
証明は，位置軌道の分解，各軌道が交代群を含むこと，全軌道の純粋工具，軌道偶奇の正規化，向きの自由度の決定，の五つの補題からなる．
本稿は，これまでプロジェクトの wiki（第3章・第5章）に記述されていた証明を，Level 2 の文書と同じ形式に移したものであり，新しい主張は加えていない．
\end{abstract}

\tableofcontents

\section{設定}

\subsection{位置，向き，状態}

Level $n$ の格子の一辺は $3^n$ である．各位置は
\[
p=\sum_{j=0}^{n-1}3^jd_j,\qquad d_j\in\{-1,0,1\}^3
\]
と一意に書ける．ベクトル $d$ の零成分の数を $z(d)$ とし，メンガー位置集合を
\[
P_n=\{\,p:\ z(d_j)\le1\ (0\le j<n)\,\}
\]
と定める．$|P_n|=20^n$ であり，$|P_3|=8000$ である．Level 3 では $p=9B+3b+o$ と書く．

立方体の回転群を $C$ と書く（$|C|=24$）．セルを完成位置で識別し，状態を $(\sigma,\rho)$ で表す．
$\sigma\colon P_3\to P_3$ はセルの現在位置，$\rho\colon P_3\to C$ は向きである．完成状態は $\sigma=\mathrm{id}$，$\rho\equiv I$ である．

\subsection{原子手と合法性}

回転対象 $D\subseteq P_n$，回転軸 $a$，軸上の基準点 $c$，角度 $\theta\in\{90^\circ,-90^\circ,180^\circ\}$ の組を原子手という．
$D$ 上で $p\mapsto c+R_{a,\theta}(p-c)$ を作用させ，向きには左から $R_{a,\theta}$ を掛ける．

原子手は，各セルを閉じた単位立方体とみたとき，回転の途中で回転対象と固定セルが正の体積で重ならない場合に合法であるという．境界での接触は許す．

Level 3 の候補対象は5,925個で，そのうち合法なものは1,197個である（表\ref{tab:families}）．
各対象に三つの角度があるので，合法原子手は $1197\times3=3591$ 個である．この集合を $S_3$，生成される群を $G_3$ と書く．
本稿で「到達可能」とは，完成状態から $S_3$ の語で得られることをいう．

\begin{table}[ht]
\centering
\caption{Level 3 の回転対象．}\label{tab:families}
\begin{tabular}{lrrl}
\toprule
族 & 候補 & 合法 & 意味 \\
\midrule
フレーム（スケール1） & 81 & 81 & 一セル厚の層 \\
フレーム（スケール3） & 27 & 27 & 三セル厚の層 \\
フレーム（スケール9） & 9 & 9 & 九セル厚の層 \\
拡張（深さ1） & 12 & 12 & 最外の辺ブロック \\
薄板（深さ1.5） & 36 & 36 & その薄板 \\
拡張（深さ2） & 240 & 72 & 中ブロック \\
薄板（深さ2.5） & 720 & 216 & 中ブロックの薄板 \\
拡張（深さ3） & 4800 & 744 & 一セル回転 \\
\midrule
合計 & 5925 & 1197 & \\
\bottomrule
\end{tabular}
\end{table}

\subsection{記法}

二つの操作 $A,B$ の交換子を $[A,B]=ABA^{-1}B^{-1}$，$T$ の $U$ による共役を $UTU^{-1}$ と書く．

\begin{lemma}[交換子の偶奇]\label{lem:comm}
有限集合上の任意の置換 $A,B$ に対して $[A,B]$ は偶置換である．
\end{lemma}
\begin{proof}
$\sgn([A,B])=\sgn(A)^2\sgn(B)^2=1$．
\end{proof}

\begin{lemma}[共役による支持の運搬]\label{lem:conj}
$T$ が点集合 $S$ の外を固定するなら，$UTU^{-1}$ は $U(S)$ の外を固定する．
\end{lemma}
\begin{proof}
$x\notin U(S)$ なら $U^{-1}(x)\notin S$ なので $T(U^{-1}(x))=U^{-1}(x)$，したがって $UTU^{-1}(x)=x$．
\end{proof}

\section{主定理}

\begin{theorem}[Level 3 の構成的可解性]\label{thm:main}
完成状態から $G_3$ の作用で到達可能な任意の状態に対し，\code{level3-slice-reduction} は，有限個の合法原子手からなる完成状態への語を構成する．
ただし，有限探索を実行する計算機がメモリ不足などの資源障害を起こさないものとする．
\end{theorem}

これは「有限手の解が存在する」だけでなく，実装がその手列を構成するという主張である．

\paragraph{適用範囲．}
定理は $S_3$ が生成する $G_3$ についてのものである．
候補生成の外にある一層の正方形断面の回転（例：$z=0$ 面の四つのセル $(\pm5,\pm5,0)$ を原点を通る Z 軸まわりに回す操作）は合法性の判定を通るが，$S_3$ には含まれない．
それらを加えた群に対する可解性は，この定理からは得られない．

\section{証明の設計}

状態 $(\sigma,\rho)$ を解く問題を，(1) 位置軌道，(2) 軌道ごとの偶奇，(3) 軌道内部の位置，(4) 位置を固定した後の向き，の四段階に分ける．

\begin{lemma}[軌道]\label{lem:orbit}
$P_3$ を頂点集合とし，ある $s\in S_3$ について $s(x)\ne x$ となるとき $x$ と $s(x)$ を辺で結んだ無向グラフの連結成分は，$G_3$ の位置軌道と一致する．
\end{lemma}
\begin{proof}
同じ連結成分にある二点は，各辺の生成元またはその逆元を順に施して移り合う．
逆に，群の元は生成元と逆元の積なので，移り合う二点の間には対応する道がある．
\end{proof}

軌道を $\Omega_1,\dots,\Omega_r$ とし，偶奇写像を
\[
\varepsilon\colon G_3\to\F_2^r,\qquad
\varepsilon(g)_i=\begin{cases}0&(g|_{\Omega_i}\text{ が偶})\\1&(g|_{\Omega_i}\text{ が奇})\end{cases}
\]
と定める．符号の乗法性から $\varepsilon(gh)=\varepsilon(g)+\varepsilon(h)$ である．

\begin{definition}[純粋位置工具]
軌道 $\Omega$ の純粋位置工具とは，$T\in G_3$ で次を満たすものをいう．
\begin{enumerate}
\item $\Omega$ の外では位置も向きも変えない．
\item $\Omega$ の中でも，位置を固定するセルの向きを変えない．
\item $\Omega$ の少なくとも三点を非自明に置換する．
\end{enumerate}
\end{definition}

\begin{lemma}[3-cycle への分解]\label{lem:threecycles}
任意の有限偶置換は3-cycleの積で表せる．
\end{lemma}
\begin{proof}
置換を巡回置換に分解する．長さ3以上の巡回置換は，先頭の点を共有する3-cycleで順に短くできる．
残る互換は偶数個であり，二つずつまとめて高々二個の3-cycleで表せる．
\end{proof}

\section{位置軌道}

\begin{proposition}[軌道分解]
3,591個の原子手の位置への作用から補題\ref{lem:orbit}のグラフを作ると，8,000個の位置は164個の軌道に分かれる（表\ref{tab:orbitsizes}）．
\end{proposition}

\begin{table}[ht]
\centering
\caption{軌道の大きさ．}\label{tab:orbitsizes}
\begin{tabular}{rrr}
\toprule
軌道の大きさ & 軌道数 & セル数 \\
\midrule
8 & 4 & 32 \\
24 & 112 & 2688 \\
96 & 45 & 4320 \\
192 & 1 & 192 \\
384 & 2 & 768 \\
\midrule
合計 & 164 & 8000 \\
\bottomrule
\end{tabular}
\end{table}

三つの桁 $B,b,o$ のそれぞれを C（零成分なし）または E（零成分一つ）に分類し，E が二つ以上あるときは零成分の軸が一致するかどうかも記録する．
例えば EEE/$Bbo$ は三桁の零軸がすべて一致し，EEE/$Bb|o$ は $B,b$ が一致して $o$ が異なることを表す．
この15クラスは座標回転で保たれる．クラスごとの軌道数は表\ref{tab:classes}のとおりである．
同じクラスの中でも軌道は分かれるので，以下の証明の単位は軌道である．164軌道の一覧は wiki の付録A（\code{wiki/Level-3-Orbit-Catalogue.md}）にある．

\begin{table}[ht]
\centering
\caption{15クラスと軌道．}\label{tab:classes}
\begin{tabular}{lrrl}
\toprule
クラス & セル数 & 軌道数 & 軌道の大きさ \\
\midrule
CCC & 512 & 24 & $8\times4,\ 24\times20$ \\
CCE & 768 & 32 & $24\times32$ \\
CEC & 768 & 26 & $24\times24,\ 96\times2$ \\
CEE/$bo$ & 384 & 13 & $24\times12,\ 96\times1$ \\
CEE/$b|o$ & 768 & 26 & $24\times24,\ 96\times2$ \\
ECC & 768 & 8 & $96\times8$ \\
ECE/$Bo$ & 384 & 4 & $96\times4$ \\
ECE/$B|o$ & 768 & 8 & $96\times8$ \\
EEC/$Bb$ & 384 & 1 & $384\times1$ \\
EEC/$B|b$ & 768 & 8 & $96\times8$ \\
EEE/$Bbo$ & 192 & 1 & $192\times1$ \\
EEE/$Bb|o$ & 384 & 1 & $384\times1$ \\
EEE/$Bo|b$ & 384 & 4 & $96\times4$ \\
EEE/$B|bo$ & 384 & 4 & $96\times4$ \\
EEE/$B|b|o$ & 384 & 4 & $96\times4$ \\
\bottomrule
\end{tabular}
\end{table}

\section{各軌道が交代群を含むこと}

軌道 $\Omega$ 上に $G_3$ が誘導する置換群を $G_3^\Omega$ と書く．

\begin{lemma}\label{lem:alt}
164個の軌道すべてについて $\Alt(\Omega)\le G_3^\Omega$ である．
\end{lemma}

\begin{proof}[計算機による証明]
大きさ8の四つの軌道では，生成される群を閉包まで列挙し，位数が $8!=40320$ であること，すなわち $\Sym(8)$ と一致することを確認した．

大きさ24以上の各軌道では，次の三点を証明書として検査した．
\begin{enumerate}
\item 作用が原始的である．すなわち生成元が保つ非自明なブロック系が存在しない．
任意の二点を同じブロックに入れた最小の合同関係を閉包し，常に全体分割しか得られないことで判定する．
\item ある生成元の語 $g$ とその冪 $g^K$ について，$g^K$ の支持がちょうど素数 $p$ 個の点からなる一つの巡回置換である．これを直接検査する．
\item $p\le|\Omega|-3$ である．
\end{enumerate}
原始的な有限置換群が，$p\le n-3$ を満たす素数長 $p$ の巡回置換を含めば交代群を含む（Jordan の定理）ので，結論が従う．
乱数は証明書の候補 $g$ を探すためにだけ使い，最終的な判定は得られた具体的な置換の検査による．
\end{proof}

\section{純粋工具}

\begin{lemma}\label{lem:tools}
各軌道 $\Omega_i$ に対し，$\Omega_i$ の外の全セルの位置と向きを固定する非自明な偶置換の工具 $T_i\in G_3$ が存在する．
\end{lemma}

工具の支持の大きさの分布は，3点が133軌道，4点が17軌道，6点が14軌道である．
117個の工具は，支持の小さい二つの交換子の組み合わせから得た．
残る47個は，工具 $T$ とその共役 $gTg^{-1}$ の交換子によって軌道の外の共通の作用を打ち消す方法（共役分離）で得た．
各工具の語は16--84原子手（中央値16）である．

\begin{proof}[計算機による証明]
保存された工具データを読み込むとき，全8,000位置について次を再検査する．
\begin{itemize}
\item 語のすべての原子手が，現行の合法原子手に解決される．
\item 動く位置が宣言された軌道の中だけである．
\item 位置を固定するセルの向きが変わらない．
\item 支持の大きさが3以上である．
\end{itemize}
独立の監査の結果は，164個中164個が合格，軌道の外への作用0件，固定セルのねじれ0件，退化した工具0件である．
\end{proof}

支持が4点または6点の工具しかない軌道では，二つの共役を合成して共通の作用を打ち消し，支持3点の工具を導く．
実装は有限探索の後に，全支持上の作用を検査する．

\section{偶奇の正規化}

状態 $s\in G_3$ の164ビットの偶奇ベクトルを $\varepsilon(s)$ とする．
3,591個の原子手が与える非零ベクトルは20種類あり，その張る空間の次元は15である．

\begin{lemma}\label{lem:parity}
任意の到達可能状態の軌道偶奇は，高々15個の原子手ですべて偶に正規化できる．
\end{lemma}
\begin{proof}
$s\in\langle S_3\rangle$ なので，$\varepsilon(s)$ は $\{\varepsilon(a):a\in S_3\}$ の張る空間に属する．
次元15の基底に対して掃き出し法を行えば，$\varepsilon(s)$ を表す基底ベクトルの和が得られる．
各基底ベクトルに対応する原子手を逆向きに施せば $0$ になる．基底の元は15個以下なので，必要な手数も15以下である．
\end{proof}

以後に使う純粋工具は交換子から作られ，補題\ref{lem:comm}により各軌道で偶なので，正規化した偶奇を壊さない．

\section{位置の解法}

各軌道について，工具を共役して，未完成のセルを完成位置へ運ぶ．
準備の語は，順序付き二点 $(x,y)\in\Omega^2$（$x\ne y$）の有限グラフを幅優先探索して求める．

支持が4点または6点の工具だけでは，最後の少数のセルで貪欲法が循環し得る．
そこで，残りが32セル以下になったら，残りの置換を補題\ref{lem:threecycles}により3-cycleへ厳密に分解する．
必要な順序付き三点上の3-cycleは，目的の三点の側と既知の工具の三点の側からの双方向の幅優先探索で構成する．

\begin{lemma}\label{lem:position}
偶奇の正規化の後，位置の解法は有限回で $\sigma=\mathrm{id}$ に到達する．
\end{lemma}
\begin{proof}
各軌道の残りの置換は偶であり，3-cycleの有限積である．
補題\ref{lem:alt}と補題\ref{lem:tools}により，必要な3-cycleの共役が存在する．
順序付き二点・三点の状態の数は有限で，探索に深さや頂点数の打ち切りがないので，存在する準備の語は必ず見つかる．
工具は軌道の外を固定するので，他の軌道を壊さない．164軌道を順に処理すれば全位置が完成する．
\end{proof}

\section{向き}

各軌道の $\Omega\times C$ 上のオートマトン（一つのセルの位置と向きの遷移）を全探索した結果，
148軌道（7,224セル）では完成位置で可能な向きが一つだけであり，位置が決まれば向きも決まる．
非自明なのは16軌道（776セル）である（表\ref{tab:freedom}）．

\begin{table}[ht]
\centering
\caption{向きの自由度が非自明な軌道．}\label{tab:freedom}
\begin{tabular}{crrl}
\toprule
群 & 軌道数 & セル数 & 処理 \\
\midrule
$C_3$ & 4 & 32 & 位置を固定するねじり工具 \\
$C_4$ & 11 & 552 & 主に合法な一セル回転 \\
$D_4$ & 1 & 192 & 一セル回転と専用のねじり工具 \\
\bottomrule
\end{tabular}
\end{table}

12軌道には合法な一セル回転がある．
大きさ8の四つの CCC 軌道は $C_3$ の自由度を持つが一セル回転がなく，工具どうしの向きへの作用の差からねじり工具を作る．
192セルの EEE/$Bbo$ 軌道は $D_4$ の自由度を持ち，一セル回転だけではすべての残差を直せないため，13個の専用の語を保存している．
監査では13個すべてが，位置を固定する非自明な向きの操作である．

\begin{lemma}\label{lem:orientation}
位置が完成した後の任意の到達可能な向きの状態は，有限回で完成する．
\end{lemma}
\begin{proof}
各軌道で到達可能な向きは，有限オートマトンにより完全に列挙されている．
合法な一セル回転，二つの純粋工具の差，$[R,T]$ 型の交換子，および $D_4$ 専用の語が，必要な部分群を生成する．
各セルの完成の向きまでの最短距離の和を $\Phi$ とする．
単一のセルで $\Phi$ を下げられない場合は，順序付き二セルの状態を有限探索し，$\Phi$ を真に減らす操作を得る．
$\Phi$ は非負整数で毎回減るので，有限回で $0$ になる．
\end{proof}

\section{主定理の証明}

\begin{proof}[定理\ref{thm:main}の証明]
入力の状態は到達可能なので，補題\ref{lem:parity}により軌道偶奇を正規化できる．
補題\ref{lem:position}により全セルの位置を有限回で完成できる．
補題\ref{lem:orientation}により，位置を固定したまま全セルの向きを有限回で完成できる．
使う語は現行の合法原子手だけからなる．したがって構成された有限の語は入力を完成状態へ移す．
\end{proof}

\section{健全性の境界検査}

探索の本体は疎な整数シミュレータを使うが，成功の判定はそれに依存しない．
入力の状態を複製し，出力を一手ずつ実エンジンに施す．各手の直前に合法性を検査し，最後に全8,000セルの位置と向きの完全な一致を要求する．
一手でも違法であるか，一セルでも一致しなければ，結果は失敗となる．
したがって，仮に探索のシミュレータに実装の差があっても，誤った成功は境界で拒否される．
さらに，疎なシミュレータと実エンジンは，6試行 $\times$ 40合法原子手 $\times$ 8,000セルについて，位置と向きの差分検査を通過している．

\section{手数}

七つの検証例では，解は111,830--424,519手，求解時間は112.1--284.2秒であった．これは最短手数ではない．

一手を合法原子手一回と数えると，長さ $d$ 以下の語の数は高々
\[
1+M+\cdots+M^d=\frac{M^{d+1}-1}{M-1},\qquad M=3591
\]
である．一方，位置の部分だけでも
\[
\prod_{i=1}^{164}\frac{|\Omega_i|!}{2}
\]
個以上の状態があり，その常用対数は約 $11394.834$ である．
したがって，最悪の状態を解くのに必要な手数（God's number）は $3206$ 以上である．
現在の解は最適でない解法の出力であり，その手数は God's number の評価としては意味を持たない．

\section{検証と再現}

\begin{itemize}
\item \code{npm run verify:l3}：保存済みの純粋工具の監査（\code{research/scratch/l3-audit-tooldata.ts}），
疎なシミュレータと実エンジンの比較（\code{l3-sim-vs-engine.ts}），軌道の置換群の証明書の検査（\code{l3-orbit-groups.ts}）．
\item \code{npm run bench -- --algorithm=level3-slice-reduction --level=3 --seeds=6 --length=100}：一回の求解の再現（数分かかる）．
\item 工具データ：\code{packages/solver-core/src/algorithms/level3SliceReductionToolData.ts}．
\item 実装：\code{packages/solver-core/src/algorithms/level3SliceReductionSolver.ts}．
\end{itemize}

\section{Level 2 の文書との関係}

Level 2 の文書（\code{level2-completeness.tex}）は，到達可能な状態の集合を不変量で完全に特徴づけ，群の位数を決定している．
本稿の Level 3 の定理は可解性だけを主張しており，$G_3$ の位数や，到達可能性の不変量による特徴づけは与えていない．
Level 2 と同じ方法（ねじれ和・反転偶奇の不変量と，三重推移性の全探索）を Level 3 に適用すれば，同じ形の特徴づけが得られると期待されるが，これは未実施である．

\end{document}
